\chapter{INTRODUCTION}

\section{Background and Motivation}

Cough is an important physiological signal that provides information about respiratory health. Continuous cough monitoring can support early detection of respiratory diseases, help in treatment follow-up, and allow long-term patient observation. However, automatic cough detection in daily-life conditions is difficult due to background noise, motion artifacts, and privacy concerns with audio recordings~\cite{dixon2025}.

Wearable sensors offer an alternative to microphone-only systems. Body-coupled audio, accelerometers, and stretch sensors can capture cough events while being less affected by environmental noise~\cite{elfaramawy2017}. Still, these signals are noisy and depend on the user's activity. Both signal processing and data-driven methods are needed for reliable detection.

This project uses a wearable system with four sensors to detect cough events and classify the user's activity during coughing. In this phase, deep learning models are developed to learn patterns directly from the sensor signals, reducing the need for manual feature engineering.

\section{Previous Work}

In the first phase of this project (EE491), a classical machine learning pipeline was developed for the same task. Time-domain and frequency-domain features were extracted from short signal windows (200\,ms duration, 50\,ms hop), and classifiers such as Random Forest and XGBoost were trained using GroupKFold cross-validation with 5 folds. The dataset at that time included 60 recordings from a single subject.

The best cough detection model (XGBoost) achieved an event-level F1-score of 0.922 using IoU-based matching at a threshold of 0.20. Activity classification between sitting and walking reached 0.89 window-level accuracy with Random Forest.

While these results were promising, the approach required manual selection of features. The model performance was limited to what the handcrafted features could capture. Deep learning methods can potentially learn better representations directly from the raw or transformed signals.

\section{Objectives}

This work extends the project to deep learning. The main objectives are:
\begin{itemize}[nosep]
    \item To design a CNN-based model architecture suitable for multi-modal sensor fusion.
    \item To investigate the effect of input representation (raw waveform vs.\ spectrogram) on cough detection performance.
    \item To introduce improved window labeling and training strategies.
    \item To compare deep learning results with the previous classical ML approach.
\end{itemize}

\section{Contributions}

The main contributions of this work are:
\begin{itemize}[nosep]
    \item A dual-branch CNN architecture that processes audio and motion signals separately and combines them through late fusion for cough detection.
    \item A center-focused window labeling strategy that reduces boundary effects in cough labels.
    \item A spectrogram-based model variant that uses 2D convolutions on log-Mel representations of audio signals.
    \item A systematic comparison of three model versions with progressively refined design choices.
\end{itemize}
